% Generated from lessons/0011-1-euclidean-space-basic-theorems.html; edit the HTML source, then rerun the converter.
\section{Euclid 空间上的基本定理}
\lessonmark{Euclid 空间上的基本定理}

这一页按教材一整节来复习：距离与点列极限、邻域与开闭集、聚点与闭包、开闭集运算、紧集与 Heine-Borel 定理，最后接教材课后题训练。

\subsection*{本节地图}

\begin{conceptbox}
\textbf{1. 距离与极限}
把一维绝对值距离推广为 \(R^n\) 中的范数距离。

\[
          |x-y|=\left(\sum_{i=1}^{n}(x_i-y_i)^2\right)^{1/2}
        \]
\end{conceptbox}

\begin{conceptbox}
\textbf{2. 点集语言}
用邻域定义内点、外点、边界点、孤立点、聚点、开集、闭集。
\end{conceptbox}

\begin{conceptbox}
\textbf{3. 基本定理}
Cantor 闭区域套、Bolzano-Weierstrass、Cauchy 收敛原理、Heine-Borel。
\end{conceptbox}

教材对应：《数学分析》第三版下册第十一章 §1“Euclid 空间上的基本定理”。辅助参考：\href{https://www.math.ucdavis.edu/~hunter/intro_analysis_pdf/ch13.pdf}{UC Davis Chapter 13}。

\subsection*{一、距离与点列极限}

在 \(R^n\) 中，一个点是一个向量 \(x=(x_1,\ldots,x_n)\)。两个点之间的距离定义为：

\[
      |x-y|=\left(\sum_{i=1}^{n}(x_i-y_i)^2\right)^{1/2}.
    \]

点列 \(x_k\to a\) 的意思是：距离 \(|x_k-a|\to 0\)。它等价于每个坐标都收敛：

\[
      x_k\to a \quad\Longleftrightarrow\quad x_k^{(i)}\to a_i,\quad i=1,\ldots,n.
    \]

\begin{notebox}
考试写法：看到 \(R^n\) 中点列极限，不要慌，把它翻译成“距离趋于 0”或“每个坐标分别趋于对应坐标”。
\end{notebox}

\subsection*{二、邻域与点的分类}

以 \(a\) 为中心、\(\delta>0\) 为半径的邻域是：

\[
      O(a,\delta)=\{x\in R^n: |x-a|<\delta\}.
    \]

\begin{center}
\small
\begin{tabularx}{\linewidth}{|>{\raggedright\arraybackslash}p{0.18\linewidth}|>{\raggedright\arraybackslash}p{0.42\linewidth}|>{\raggedright\arraybackslash}X|}
\hline
\textbf{概念} & \textbf{邻域判别} & \textbf{最该记住的话} \\ \hline
内点 & \(\exists\delta>0,\ O(x,\delta)\subset S\) & 周围有一整圈安全区域。 \\ \hline
外点 & \(\exists\delta>0,\ O(x,\delta)\subset S^c\) & 周围完全在集合外。 \\ \hline
边界点 & 任意邻域同时碰到 \(S\) 和 \(S^c\) & 怎么缩小都两边沾。 \\ \hline
孤立点 & \(\exists\delta>0,\ O(x,\delta)\cap S=\{x\}\) & 它属于 \(S\)，但附近没有同伴。 \\ \hline
聚点 & \(\forall\delta>0,\ O(x,\delta)\cap(S\setminus\{x\})\ne\varnothing\) & 每个邻域都有别的 \(S\) 点。 \\ \hline
\end{tabularx}
\end{center}

\[
      x\in S' \quad\Longleftrightarrow\quad
      \exists\{x_k\}\subset S,\ x_k\ne x,\ x_k\to x.
    \]

这里 \(S'\) 表示 \(S\) 的全体聚点。教材里特别有用的结论是：用点列可以判别聚点。

\subsection*{三、开集、闭集、闭包}

\begin{conceptbox}
\textbf{开集}
\(S\) 中每个点都是内点。

\[
          \forall x\in S,\ \exists\delta>0,\ O(x,\delta)\subset S.
        \]
\end{conceptbox}

\begin{conceptbox}
\textbf{闭集}
\(S\) 包含它的全部聚点。

\[
          S'\subset S.
        \]
\end{conceptbox}

\begin{conceptbox}
\textbf{闭包}
\(S\) 加上全部聚点。

\[
          \overline S=S\cup S'.
        \]
\end{conceptbox}

开闭集的运算是期末常考证明素材：

\begin{center}
\small
\begin{tabularx}{\linewidth}{|>{\raggedright\arraybackslash}X|>{\raggedright\arraybackslash}X|}
\hline
\textbf{结论} & \textbf{考试理由} \\ \hline
任意多个开集的并是开集 & 点落在某一个开集中，于是继承那个开集给出的邻域。 \\ \hline
任意多个闭集的交是闭集 & 用补集和 De Morgan 公式转化为开集的并。 \\ \hline
有限多个开集的交是开集 & 每个开集给一个半径，取这些半径的最小值。 \\ \hline
有限多个闭集的并是闭集 & 用补集和 De Morgan 公式转化为有限开集的交。 \\ \hline
\end{tabularx}
\end{center}

\begin{notebox}
无限交开集不一定开；无限并闭集不一定闭。反例常用 \(\bigcap_{n=1}^{\infty}(-1/n,1/n)=\{0\}\) 和 \(\bigcup_{n=1}^{\infty}[1/n,1]=(0,1]\)。
\end{notebox}

\subsection*{四、紧集与四个基本定理}

教材用开覆盖定义紧集：若 \(S\) 的任意开覆盖都能取出有限子覆盖，则 \(S\) 是紧集。

\[
      S\subset \bigcup_{\alpha}U_\alpha
      \quad\Rightarrow\quad
      S\subset \bigcup_{j=1}^{m}U_{\alpha_j}.
    \]

在 \(R^n\) 中，Heine-Borel 定理把抽象的开覆盖条件变成非常好用的判别：

\[
      S\subset R^n \text{ compact}
      \quad\Longleftrightarrow\quad
      S \text{ bounded and closed}.
    \]

教材还把以下命题放在同一组基本定理里，它们在 Euclid 空间中相互等价：

\begin{conceptbox}
\textbf{Bolzano-Weierstrass}
有界无限点集至少有一个聚点。考试里常用于“从有界点列取收敛子列”。
\end{conceptbox}

\begin{conceptbox}
\textbf{Cauchy 收敛原理}
\(R^n\) 中点列收敛，当且仅当它是 Cauchy 点列。
\end{conceptbox}

\begin{conceptbox}
\textbf{Cantor 闭区域套}
嵌套闭区域且直径趋于 0 时，交集只剩唯一一点。
\end{conceptbox}

\begin{conceptbox}
\textbf{Heine-Borel}
\(R^n\) 中紧集等价于有界闭集。
\end{conceptbox}

\subsection*{五、即时判断}

\begin{quickcheck}
若 \(S'\) 表示 \(S\) 的聚点集，则 \(S\cup S'\) 是什么？

\tcblower
\textbf{答案：}闭包。闭包就是集合加上所有聚点，所以 \(\overline S=S\cup S'\)。
\end{quickcheck}

\begin{quickcheck}
在 \(R\) 中，集合 \([0,1]\) 属于哪一类？

\tcblower
\textbf{答案：}闭非开。\([0,1]\) 含有所有聚点，所以闭；端点 \(0,1\) 不是内点，所以不开。
\end{quickcheck}

\begin{quickcheck}
在 \(R^n\) 中，一个集合有界且闭。按 Heine-Borel 定理，它是什么？

\tcblower
\textbf{答案：}紧集。在 \(R^n\) 中，Heine-Borel 定理说有界闭集等价于紧集。
\end{quickcheck}

\subsection*{六、课后题训练}

下面按教材第 106 页习题 1-9 整理为考试训练卡。数学式保留，表述做了复习化整理；写作业时请以教材原题为准。

\begin{exercisebox}
\textbf{习题 1\badge{距离公理}}

证明 \(R^n\) 中距离满足正定性、对称性和三角不等式。

\begin{hintbox}
正定性和对称性直接由平方和得到。三角不等式可用 Cauchy-Schwarz 不等式：

\[
          \left(\sum_{i=1}^{n}a_i b_i\right)^2
          \le \left(\sum_{i=1}^{n}a_i^2\right)\left(\sum_{i=1}^{n}b_i^2\right).
        \]
\end{hintbox}
\end{exercisebox}

\begin{exercisebox}
\textbf{习题 2\badge{极限唯一}}

证明：若 \(R^n\) 中点列 \(\{x_k\}\) 收敛，则极限唯一。

\begin{hintbox}
假设 \(x_k\to a\) 且 \(x_k\to b\)。用三角不等式：

\[
          |a-b|\le |a-x_k|+|x_k-b|\to 0.
        \]

所以 \(|a-b|=0\)，即 \(a=b\)。
\end{hintbox}
\end{exercisebox}

\begin{exercisebox}
\textbf{习题 3\badge{极限运算}}

设 \(x_k\to x\)，\(y_k\to y\)。证明对任意实数 \(\alpha,\beta\)，有

\[
        \lim_{k\to\infty}(\alpha x_k+\beta y_k)=\alpha x+\beta y.
      \]

\begin{hintbox}
估计距离：

\[
          |(\alpha x_k+\beta y_k)-(\alpha x+\beta y)|
          \le |\alpha|\,|x_k-x|+|\beta|\,|y_k-y|.
        \]
\end{hintbox}
\end{exercisebox}

\begin{exercisebox}
\textbf{习题 4\badge{内部边界闭包}}

求以下 \(R^2\) 子集的内部、边界与闭包：

\[
        \begin{aligned}
        &(1)\ S=\{(x,y):x>0,\ y\ne0\},\\
        &(2)\ S=\{(x,y):0<x^2+y^2\le1\},\\
        &(3)\ S=\{(x,y):0<x\le1,\ y=\sin(1/x)\}.
        \end{aligned}
      \]

\begin{hintbox}
(1) 把第一、四象限去掉 \(x\) 轴来想；边界包括 \(y=0,x\ge0\) 和 \(x=0\)。

(2) 是闭圆盘去掉原点；原点会进入边界和闭包。

(3) 是振荡曲线；当 \(x\to0^+\) 时，极限点填满线段 \(\{0\}\times[-1,1]\)。
\end{hintbox}
\end{exercisebox}

\begin{exercisebox}
\textbf{习题 5\badge{全部聚点}}

求以下点集的全部聚点：

\[
        \begin{aligned}
        &(1)\ S=\{(-1)^k\,k/(k+1):k=1,2,\ldots\},\\
        &(2)\ S=\{(\cos(2k\pi/5),\sin(2k\pi/5)):k=1,2,\ldots\},\\
        &(3)\ S=\{(x,y):(x^2+y^2)(y^2-x^2+1)\le0\}.
        \end{aligned}
      \]

\begin{hintbox}
(1) 分奇偶子列，两个聚点是 \(-1\) 和 \(1\)。

(2) 只有 5 个不同点，有限集没有聚点。

(3) 因为 \(x^2+y^2\ge0\)，条件等价于 \(y^2-x^2+1\le0\) 或 \((0,0)\) 的特殊情况；这是闭集，先画出区域。
\end{hintbox}
\end{exercisebox}

\begin{exercisebox}
\textbf{习题 6\badge{聚点序列判别}}

证明：\(x\) 是点集 \(S\subset R^n\) 的聚点，当且仅当存在 \(S\) 中点列 \(\{x_k\}\)，满足 \(x_k\ne x\)，且 \(x_k\to x\)。

\begin{hintbox}
必要性：对每个 \(k\)，在 \(O(x,1/k)\) 中选一个 \(x_k\in S\setminus\{x\}\)。充分性：若 \(x_k\to x\)，任意邻域最终都会包含 \(x_k\)。
\end{hintbox}
\end{exercisebox}

\begin{exercisebox}
\textbf{习题 7\badge{开闭反例}}

若 \(U\subset R^2\) 是开集，是否 \(U\) 的每个点都是 \(U\) 的聚点？若 \(F\subset R^2\) 是闭集，又如何？

\begin{hintbox}
非空开集中的每个点都是聚点，因为每个邻域里能取到不同于该点的点。闭集不一定，例如单点集 \(\{(0,0)\}\) 是闭集，但它的点不是聚点。
\end{hintbox}
\end{exercisebox}

\begin{exercisebox}
\textbf{习题 8\badge{内点集开}}

证明 \(S\subset R^n\) 的所有内点组成的点集 \(S^\circ\) 是开集。

\begin{hintbox}
若 \(x\in S^\circ\)，存在 \(O(x,r)\subset S\)。取 \(0<\rho<r\)，证明 \(O(x,\rho)\) 中的每个点也还是 \(S\) 的内点。
\end{hintbox}
\end{exercisebox}

\begin{exercisebox}
\textbf{习题 9\badge{闭包闭集}}

证明 \(S\subset R^n\) 的闭包 \(\overline S=S\cup S'\) 是闭集。

\begin{hintbox}
证明 \(\overline S\) 含有自己的全部聚点。若 \(x\) 是 \(\overline S\) 的聚点，则任意邻域碰到 \(\overline S\)，再由闭包点的定义推出任意邻域碰到 \(S\)，所以 \(x\in\overline S\)。
\end{hintbox}
\end{exercisebox}

\subsection*{七、期末写法模板}

\begin{conceptbox}
\textbf{证明开集}
任取 \(x\in S\)，构造 \(\delta>0\)，证明 \(O(x,\delta)\subset S\)。
\end{conceptbox}

\begin{conceptbox}
\textbf{证明闭集}
用 \(S'\subset S\)，或取 \(x_k\in S,x_k\to x\)，证明 \(x\in S\)。
\end{conceptbox}

\begin{conceptbox}
\textbf{证明聚点}
任取 \(\delta>0\)，在 \(O(x,\delta)\cap(S\setminus\{x\})\) 中构造点。
\end{conceptbox}

\begin{conceptbox}
\textbf{证明紧集}
在 \(R^n\) 中优先用 Heine-Borel：证明有界且闭。
\end{conceptbox}

下一节：第十一章 §2：多元连续函数。补充练习：用球邻域判断开闭，聚点、闭包、边界。路线图：教材路线图。
